%======================================================================
%	Vorlage
%======================================================================
%	$Id$
%	Matthias Kupfer
%======================================================================
%	Documentclass
%======================================================================
\documentclass[
% 	draft,			% Entwurfsmodus: Bilder als Rahmen,
				% Überlängen werden deutlich markiert
%	10pt,			% 
	11pt,			% KOMA default
%	12pt,			% 
	a4paper,		% DIN A4
	oneside,		% Zweiseitig
	german,			%
%----------------------------------------------------------------------
% Die folgenden Befehle stammen aus dem KOMA Paket
	headsepline,		% Linie unter der Kopfzeile	
%	headnosepline,		%
%	foodsepline,		% Linie über Fussnote	
	footsepline=false,	%
	automark,		% Kolumnentitel lebendig
%	bigheadings,		% default: Überschriften gross setzen
	headings=normal,	% Überschriften normal setzen
%	smallheadings,		% Überschriften eher klein setzen
%	pointlessnumbers,	% Keinen Punkt hinter die letzte Zahl 
				% eines Kapitels (auch bei Anhang)
%	chapterprefix,		% Kapiteluberschriften "Kapitel"
%	appendixprefix,		% Anhang
%	openright,		% Kapitel auf der rechten (ungeraden) Seite 
				% anfangen lassen
	openany,		% 
%	cleardoublestandard,	%
	cleardoublepage=plain,	%
%	cleardoubleempty,	%
	abstracton,		%
	index=totoc,		% Index soll im Inhaltsverzeichnis auftauchen
	listof=totoc,		%
	bibliography=totoc,	%
%	parskip, 		% parskip-, parskip*, parskip+
% 	halfparskip, 		% halfparskip-, halfparskip* und halfparskip+
%	DIVclassic,
 	BCOR8mm,		%%?? Bindungskorrektur: 
				% BCOR<Breite des Bindungsverlustes> 
]{scrreprt} 

%======================================================================
%	Bearbeitung sowohl mit LaTeX als auch mit pdfLaTeX ermoeglichen
%======================================================================
%\newif\ifpdf 
	%\ifx\pdfoutput\undefined 
	%\pdffalse 			% we are not running pdflatex 
%\else 
	%\pdfoutput=1 			% we are running pdflatex 
	%\pdftrue 
%\fi


%======================================================================
%	Verwendete Pakete
%======================================================================

%\usepackage{babel}		% Sprachen

%% Latex mit deutschen Umlauten:
%% http://www.cs.albany.edu/~herrmann/latex_umlaute/
\usepackage[T1]{fontenc}	% EC-Schriften verwenden (vs. DC) da 8-Bit
				% EC-Schriften als T1-kodierten CM-Schriften
				% European/Ext.-Computer-Modern-(EC)-Schriften
				% Umlaute, Anführungszeichen ...
				% => Umlauten koennen richtig getrennt werden
				% FAQ 5.3.2

%\usepackage{ucs}		% Eingabe von ü,ö,ä,ß mit UTF-8 erlaubt
\usepackage[utf8]{inputenc}	% (auch in -Dateien)
				% ! Zeichensatz der Dateien als UTF-8 !
								
\usepackage{ae,aecompl}		% virtuelle-CM-Fonts
				% da EC nicht als PostScript-(Type-1) verfuegbar
				% => keine echten Umlaute im PDF-Dokumen 
				%(Problem bei Suche)
				% By loading the ae package (\usepackage{ae}), 
				% you loose some characters as mentioned in 
				% README. 
				% The package aecompl by Denis Roegel restores
				% these characters which are taken from the ec 
				% fonts. If you use pdftex, you will get these 
				% characters as bitmaps, but this might be 
				% better than not having them at all.
\usepackage{bbold}
\usepackage{amsmath}
\usepackage[hang]{caption}
%\usepackage{times, mathptm}	% TimesNewRoman Schrift (Acrobat Reader Fonts), 
				% dazu braucht man auch den entsprechenden 
				% Zeichensatz für den Math-Mode
%\usepackage{pslatex}		% ? mathematische Formeln mit Standard 
				% Postscript Fonts gesetzt
			% Paket     Roman         Serifenlos  Typewriter
%\usepackage{times}	% -----------------------------------------------
			% times     Times         Helvetica   Courier
			% palatino  Palatino      Helvetica   Courier
			% newcent   NewCenturySch AvantGarde  Courier
			% bookman   Bookman       AvantGarde  Courier
			% Diese Schriften sind die Standard-PostScript-Schriften
			% und in jedem Drucker verfügbar

\usepackage{
%	german,			% Deutsche Trennungen (ALTE Rechtschreibung), 
				% Anführungsstriche und mehr 
	ngerman,		% Deutsche Trennungen (NEUE Rechtschreibung), 
				% Anführungsstriche und mehr
				%
%	acronym,		% Verwaltung von Abkuerzungen
	bibgerm,		% Deutsche Bibliographie 
	calc,			% Erweiterung der arithmetischen Funktionen in 
				% LaTeX
				% wird verwendet um Titelseite zu zentrieren
	color,			% im Laufenden Text einfach mit \color{Farbe) zwischen den 
				% Farben umschalten, wobei Farbe einfach 
				% durch z.B. red, blue, black etc. ersetzt wird
				% \textcolor{farbe){Text)
%	epigraph,		% Zitat am Kapitelanfang
%	fancyhdr,		% Kopf- und Fußzeilen von Dokumenten frei 
				% gestalten
	fancybox,		% shadowbox, doublebox, ovalbox, Ovalbox 
	fancyvrb,		% verbatim Erweiterung:
	float,			% Positionierung von Gleitobjekten genau an der Stelle, wo man
				% 'figure'- oder 'table'-Umgebung die 
				% Positionierung [H] gesetzt werden
%	glosstex,		% Glossar und Abkürzungsverzeichnis
	mdwlist,		% compact list: itemize* ..
	scrdate,		% \todaysname 
	scrtime,		% \thistime
	scrpage2,		% Kopf- und Fußzeilen flexibel gestalten
				%
%	moreverb,		% verbatim-ähnlich: boxedverbatim, listing
%	verbatim,		% Darstellung von "Text, wie er eingegeben wird"
				%
%	lscape,			% Erstellt eine um 90% gedrehte *neue* Seite
%	textcomp,		% Sonderzeichen
%	booktabs,		% Tabellenlinien
%	longtable,		% Tabellen > 1 Seite
%	supertabular,		% Tabellen > 1 Seite
	tabularx,		% Blocksatzspalten
%	ltxtable,		% tabularx + longtable
%	multicol,		% mehrspaltige Zeilen
%	varioref,		% einheitliche Verweise
%	endnotes,		% Fussnoten -> Endnoten
%	rotating,		% sidewaystable und sidewaysfigure
%	natbib,			% Bibliographie ohne Klammer etc.
%	marvosym,		% Euro etc.
}

%\usepackage{pstricks}
%\usepackage{listings}
%\usepackage{wasysym}

%\usepackage[german, first, bottomafter]{draftcopy}

%\usepackage{setspace}	% Durchschuß, Zeilenabstand
%\doublespace		% doppelzeilig oder
%\onehalfspacing	% anderthalbzeilig
\renewcommand{\baselinestretch}{1.1}\normalsize

% Für schöne Darstellung von Algorithmen
%\usepackage[german, algoruled, algochapter]{algorithm2e}
%\usepackage{algorithmic}
%\usepackage[chapter]{algorithm}
%\floatname{algorithm}{Algorithmus}

%======================================================================
%	Bilder, Links
%======================================================================

\usepackage{graphicx}
\graphicspath{Bilder/}	% Angabe der Pfade, wo die Grafiken liegen; 
				% mehrere Pfade sind möglich
\usepackage{subfig}
\usepackage{thumbpdf}
\usepackage{hyperref}
%\usepackage{breakurl}


%======================================================================
%	user.sty
%======================================================================

\usepackage{user}	% Makros in user.sty 
			% \epsinc{bild.eps}{scale=1}{Bildunterschrift}
			% \missing{Beweis fehlt noch}
			% \comment{Ein Kommentar}

%======================================================================
%	Einstellungen
%======================================================================

%\typearea{11}		% Satzspiegel neu konstruieren (KOMA)
			%      10pt 11pt 12pt 
			% DIV : 8   10   12  

\pagestyle{scrheadings}	% Standart  Kopf- und Fußzeile
\setkomafont{pagehead}{\small\scshape}

% ---------------------------------------------------------------------
%\setcounter{secnumdepth}{2}
%\setcounter{chapter}{-1}
\setcounter{tocdepth}{3}

% ---------------------------------------------------------------------
%\sloppy		% weniger Worttrennungen, größere Wortabstände
\fussy			% viele Worttrennungen, "schönere" Wortabstände

% ---------------------------------------------------------------------
%\flushbottom			% Ausrichtung der Seitenenden jeweils auf 
				% gleicher Höhe

% ---------------------------------------------------------------------
%\sloppypar		% Das hier relaxt die Einstellungen zum Wortabstand 
			% extrem. Damit ragen keine Worte über den rechten 
			% Zeilenabstand hinaus. Dafür muß stärker auf
			% Wortabstand geachtet werden, der kann dann ziemlich 
			% groß werden. Man erhält aber keine Meldung mehr 
			% über underfull boxes.

%Hiermit kann man das gleiche mit weniger Holzhammer erreichen:
%\setlength{\tolerance}{2000}           % Strafpunkt für Zeilenumbruch
%\setlength{\emergencystretch}{3pt}     % Soweit dürfen einzelne Worte mehr 
					% auseinandergezogen werden
%\setlength{\hfuzz}{1pt}                % Macht den rechten Rand um bis zu 1pt 
					% flatterig.

% ---------------------------------------------------------------------
%Vermeiden einzelner Zeilen am Ende einer Seite oder oben auf einer neuen Seite
%\clubpenalty10000
%\widowpenalty10000

% ----------------------------------------------------------------------
% neue Umgebungen für verwendete Sätze und Beispiele
\newtheorem{bsp}{Beispiel}[chapter]
\newtheorem{satz}{Satz}[chapter]

% Eigene Anpassungen
% Silbentrennung
\tolerance 1414
\hbadness 1414                                                                  \emergencystretch 1.5em
\hfuzz 0.3pt
\widowpenalty=10000
\hyphenpenalty=200
\clubpenalty = 10000
\widowpenalty = 10000
\displaywidowpenalty = 10000
\vfuzz \hfuzz
\raggedbottom


%======================================================================
%	includeonly
%======================================================================

\includeonly{		% Gibt an, welche Dateien der include-Befehl 
			% tatsächlich einfuegen darf.
  metadaten		% Variablen setzen
  ,titel		% Titelseite, Zusammenfassung und Inhaltsverzeichnis
% Alle per include einzulesenden Dateien müssen hier angegeben sein!
  ,IM/Projektbericht
}

% Aufteilung der Arbeit in folgende Bestandteile in angeg. Reihenfolge
%
% Titelseite
% Bibliographische Beschreibung (Rückseite der Titelseite)
% (Aufgabenstellung)
% Danksagung
% Abstract
% Inhaltsverzeichnis
% Tabellenverzeichnis
% Abbildungsverzeichnis
% Einleitung
% Inhalt
% Zusammenfassung/Ausblick
% (Verzeichnis verwendeter Terme)
% Glossar (Abkürzungsverzeichnis)
% (Index)
% Literaturverzeichnis
% (Thesen)
% (Selbstständigkeitserklärung)
% Anhänge

% \renewcommand{\familydefault}{rob}
% \renewcommand{\sfdefault}{rob}

%======================================================================
% * *   T E X T    S T A R T S    H E R E   * * * * * * * * * * * * * *
%======================================================================

\begin{document}
\deffootnote[1em]{1em}{1em}{\textsuperscript{\thefootnotemark}}
%======================================================================
%	Metadaten
%======================================================================
%	$Id$
%	Matthias Kupfer
%======================================================================

\newcommand{\dcsubject}{Projektbericht}
% z.B. (Diplom/Studien/Haus)arbeit, Praktikumsbericht, Studie, Beleg, 
% (Pro/Haupt/Ober)seminar, Seminar usw.
\newcommand{\dctitle}{Erstellung eines interaktiven, nicht-linearen Videos im Rahmen der Vorlesung \glqq Interactive Media\grqq}
\newcommand{\dcsubtitle}{~} % Untertitel, falls erforderlich

\newcommand{\dcauthorlastname}{Conrad}
\newcommand{\dcauthorfirstname}{Steve}
\newcommand{\dcauthormatrikelnummer}{234498}
\newcommand{\dcauthoremail}{steve.conrad@s2009.tu-chemnitz.de} 
\newcommand{\dcdate}{\today} 

\newcommand{\dcplace}{Chemnitz} % Ort, kann an der TU meist so bleiben
\newcommand{\dcuni}{Technische Universität \dcplace}
\newcommand{\dcdepart}{Fakultät für Informatik} % Fakultätsangabe
\newcommand{\dcprof}{Professur Medieninformatik} % Angabe der Professur

\newcommand{\dcpruefer}{Dr. Britta Meixner}% Prüfer der Arbeit
\newcommand{\dcadvisor}{Dipl.-Inf. Robert Manthey}% Betreuer der Arbeit

\newcommand{\dckeywords}{multimedia, interactive, media, hypervideo, clickable, video, annotation, hot spots, streetview, bruehl, chemnitz}

%%======================================================================
% Einstellungen des Hyperref-Paketes
\hypersetup{%    
	pdftitle	= {\dctitle}, %
	pdfsubject	= {\dcsubject, \dcdate}, %
	pdfauthor	= {\dcauthorfirstname~\dcauthorlastname, \dcauthoremail}, %
	pdfkeywords	= {\dckeywords}, %
	pdfcreator	= {pdfTeX with Hyperref and Thumbpdf}, %
	pdfproducer	= {LaTeX, hyperref, thumbpdf}, %
	% weitere PDF-Einstellungen in hyperref.cfg
}
%%======================================================================
 % Variablen, \hypersetup, etc.

%======================================================================
%	Titelseite (Zusammenfassung und Inhaltsverzeichnis)
%======================================================================

%======================================================================
%	Titelseite 
%======================================================================
%	$Date:$
%	$Revision:$
%	Matthias Kupfer
%======================================================================

%%======================================================================
%% Schmutztitel
%%======================================================================
%\extratitle{
%	\usekomafont{disposition}\mdseries 
%	\begin{center}
%		\Huge \dcsubject\\[1.5ex]
%		\hrule
%		\vspace*{\fill}
%		\includegraphics{TUC_deutsch_einzeile_CMYK}
%	\end{center}
%}

%%======================================================================
%% Titelkopf
%%======================================================================
\titlehead{
	\vspace*{-1.5cm}
	% Schriftfamilie wie alle Überschriften, aber nicht fett
	\usekomafont{disposition}\mdseries 
	\begin{center}
		\raisebox{-1ex}{\includegraphics[scale=1.2]{TU_Chemnitz_positiv_gruen}}\\
		\hrulefill \\[1em]
		{\Large\dcdepart}\\[0.5em] 
		\dcprof
	\end{center}
	\vspace*{1.5cm}
}

%%======================================================================
%% Subjekt
%%======================================================================
\subject{\normalfont \bfseries\Huge\dcsubject}


%%======================================================================
%% Titel
%%======================================================================
\title{\normalfont \LARGE
	\dctitle
	\dcsubtitle
}

%%======================================================================
%% Autor des Dokumentes
%%======================================================================
\author{\dcauthorfirstname~\dcauthorlastname \\ Matrikelnummer: \dcauthormatrikelnummer \\ \\ Benjamin Kindt \\ Matrikelnummer: 275059 \\ \\}

	
%%======================================================================
%% Ort, Datum
%%======================================================================
\date{\dcplace, den \dcdate
}

%%======================================================================
%% Publishers
%%======================================================================
\publishers{
	{\parbox{\textwidth-8em}{
		\begin{tabbing}
			{\normalfont \bfseries Betreuer:}\quad\=\kill
			{\normalfont \bfseries Prüfer:}	\>\dcpruefer\\
			%{\normalfont \bfseries Betreuer:}	\>\dcadvisor
		\end{tabbing}	
	}}
}

%%======================================================================
%% maketitle
%%======================================================================

\maketitle

%%======================================================================
%% bibliografische Angaben
%%======================================================================
%\newpage
%\thispagestyle{empty}
%\null\vfill
%\null\vfill
%\noindent \textbf{\dcauthorlastname, \dcauthorfirstname}\\
%\dctitle\\
%\dcsubject,~\dcdepart\\
%\dcuni,~\ifcase\month\or
%  Januar\or Februar\or März\or April\or Mai\or Juni\or
%    Juli\or August\or September\or Oktober\or November\or Dezember\fi
%    ~\number\year


%%======================================================================
%% Danksagung
%%======================================================================
%\newpage
%\thispagestyle{empty}
%\null\vfil
%\begin{center}
%\usekomafont{disposition}\textbf{Danksagung}
%\vspace{-.5em}\vspace{\parsep}
%%
%% Hier steht der Text für die Danksagung
%%
%\input{Text/Danksagung}
%\end{center}
%\par\vfil\null
%\cleardoubleemptypage
  
%%======================================================================
%%      Kurzfassung / Abstract
%%======================================================================
%\def\abstractname{Abstract} 	% Wenn der Text "Zusammenfassung" erscheinen 
%				% soll, dann muß dies auskommentiert werden
%				
%\begin{abstract}
%%
%% Inhalt der Arbeit

%%
%\end{abstract}

%%======================================================================
%%      Inhaltsverzeichnis
%%======================================================================
\cleardoubleemptypage
\pagenumbering{roman}
%\pdfbookmark{Inhaltsverzeichnis}{Inhaltsverzeichnis}
\tableofcontents

%%======================================================================
%%      Abbildungsverzeichnis
%%======================================================================
%\cleardoublepage
%\markboth{Abbildungsverzeichnis}{Abbildungsverzeichnis}
%\listoffigures

%%======================================================================
%%      Tabellenverzeichnis
%%======================================================================
\cleardoublepage
\markboth{Tabellenverzeichnis}{Tabellenverzeichnis}
\listoftables

%%======================================================================
%%      Algorithmenverzeichnis
%%======================================================================
%\renewcommand{\listalgorithmname}{Algorithmenverzeichnis}
%\cleardoublepage
%\addcontentsline{toc}{chapter}{Algorithmenverzeichnis}
%\listofalgorithms

%%======================================================================
%%      Abkuerzungsverzeichnis
%%======================================================================
%\cleardoublepage
%\addcontentsline{toc}{chapter}{Abkürzungsverzeichnis}
%\markboth{Abk"urzungsverzeichnis}{Abk"urzungsverzeichnis}
%\def\listacronymname{Abk"urzungsverzeichnis}
%\printglosstex(acr)

%%======================================================================
%%      Ende
%%======================================================================
\cleardoublepage
\pagenumbering{arabic}
\setcounter{page}{1}
 % Titelseite, Zusammenfassung 
		% Inhaltsverzeichnis
		% Abkuerzungsverzeichnis
		% -> danach Seitennummerierung bei 1

%======================================================================
%	Inhalt per include
%======================================================================

%\include{Text/Danksagung}
%\include{LatexBeispiele/Einleitung}
\chapter{Einführung und Projektidee}
Im Rahmen der Vorlesung \glqq Media Interactive\grqq wurden wir vor die Aufgabe gestellt, ein interaktives, nicht-lineares Video zu erstellen. Das Thema des Videos konnten wir frei nach unseren Interessen wählen.Wir entschieden uns dazu, ein Video über den Brühl Boulevard, auf dem wir beide wohnen, zu erstellen.
\\ \\
Der Brühl war zur Zeit der Deutschen Demokratischen Republik und auch nach der Wende ein beliebtes Wohn- und Gewerbeviertel. Mit zahlreichen Geschäften, vielen Arbeitsplätzen sowie für die damalige Zeit modernen Wohnungen zog die Flaniermeile Menschen von Nah und Fern an. Leider verlor der Brühl nach der Wende durch die Konzentration des Einzelhandels in Einkaufszentren am Stadtrand sowie einer verfehlten Stadtentwicklungsstrategie maßgeblich an Attraktivität, was in Verbindung mit dem demografischen Wandel zu einem enormen Leerstand führte. Seit einigen Jahren jedoch plant die Stadt eine Reaktivierung des Brühls als innerstädtisches Wohngebiet für alle Generationen und Zielgruppen. Durch das gleichzeitig hohe private Engagement der Anwohner für eine Belebung ihres Boulevard sind erste Erfolge bereits deutlich sichtbar.
\\ \\
Eine Möglichkeit Menschen den Brühl noch weiter näherzubringen wäre eine virtuelle, videobasierte Tour über den Boulevard. Angelehnt an Google's bildbasierten Street View Dienst, könnte man unsere videobasierte Tour über den Brühl weitaus umfangreicher umsetzen: 360 Grad Videoaufnahmen, gegebenenfalls in Kombination mit Virtueller Realität und 3D-Technik könnten die Immersion in die virtuelle Umgebung erhöhen. Weiterhin sind Aufnahmen in Abhängigkeit von der Tags- und Jahreszeit oder tagesaktuelle Aufnahmen, zum Beispiel von Veranstaltungen, denkbar. Man müsste sich auch nicht auf den Boulevard an sich beschränken, sondern könnte auch Innenaufnahmen von Geschäften und Läden oder Luftaufnahmen, zum Beispiel die einer Drohne, in das Projekt integrieren. Interessant wäre zudem ein direkter Vergleich aktueller Aufnahmen mit historischen Bild- und Videoaufnahmen aus der identischen Aufnahmeposition und der gleichen Perspektive.
\\ \\
 Eine vielfältige Endgeräteunterstützung würde den Zugang zum Projekt erleichtern. Neben den klassischen Geräten wie PC, Smartphone und Tablet sind auch moderne Formen der multimedialen Endgeräte denkbar, wie zum Beispiel VR-Headsets oder Google Glass.



\chapter{Design und Konzeption}
Im Fokus unseres Players steht das Wiedergabefenster mit dem entsprechend gerade abgespielten Video. Darunter befindet drei Boxen, die Informationen zum Video bereitstellen und/oder eine Interaktion zulassen. In der ersten Box auf der linken Seite wird dem Benutzer die Blickrichtung sowie der Ort der Aufnahme in Form der Straßenbezeichnung und der geographischen Koordinaten angezeigt. Die mittlere Box ermöglicht das navigieren im interaktiven Video. Auf diese Weise kann sich der Benutzer virtuell über den Brühl Boulevard bewegen. Falls zum gerade angezeigten Video ein Point-of-Interest zur Verfügung steht, wird auf der rechten Seite eine weitere Box eingeblendet. In dieser wird die entsprechende Bezeichnung und Beschreibung des POI's sowie ein Link zu einer Website mit weiterführenden Informationen angezeigt.\\ \\
Statt die Navigationspfeile in einer Box unter dem Wiedergabefenster unterzubringen, wäre es auch denkbar gewesen, diese als Overlay über dem Video anzuzeigen. Das selbe gilt für Overlay-Informationen über entsprechend vorhandenen Points-of-Interests. \\
Weiterhin wäre es möglich gewesen, anhand der zu jedem Video verfügbaren geographischen Koordinaten eine interaktive Karte im Player einzublenden, die sowohl die aktuelle Position als auch alle anderen Aufnahmepositionen visualisiert sowie durch einfaches Klicken zusätzlich die Möglichkeit der Navigation zu diesen ermöglicht hätte.



\chapter{Ablaufplan}
In der Tabelle \ref{videoübersicht} auf der folgenden Seite sind die insgesamt 26 im Projekt verwendeten Videos aufgelistet. Zu jedem Video ist die Video-ID, der Aufnahmestandort, die Blickrichtung in Form der bekannten Himmelsrichtungen, die mit dem Video verknüften Nachbarvideos sowie, falls vorhanden, der Point-of-Interest angegeben.
\begin{table}[]
\centering
\caption{Videoübersicht}
\label{videoübersicht}
\begin{tabular}{lllll}
ID & Position                                                                    & Blickrichtung & POI                                                                              & Nachbarn                                                     \\
2  & Brühl/Georgstraße                                                           & WNW           & -                                                                         & 3, 4                                                         \\
3  & Brühl/Georgstraße                                                           & SSW           & -                                                                                & 2, 4                                                         \\
4  & Brühl/Georgstraße                                                           & NNO           & Brühltor                                                                                & 2, 3, 8                                                      \\
5  & \begin{tabular}[c]{@{}l@{}}Brühl/Untere\\ Aktienstraße\end{tabular}         & WNW           & -                                                                                & 6, 7, 8                                                      \\
6  & \begin{tabular}[c]{@{}l@{}}Brühl/Untere\\ Aktienstraße\end{tabular}         & SSW           & -                                                                                & 3, 5, 7, 8                                                   \\
7  & \begin{tabular}[c]{@{}l@{}}Brühl/Untere\\ Aktienstraße\end{tabular}         & OSO           & -                                                                                & 5, 6, 8                                                      \\
8  & \begin{tabular}[c]{@{}l@{}}Brühl/Untere\\ Aktienstraße\end{tabular}         & NNO           & -                                                                                & 5, 6, 7, 9                                                   \\
9  & Brühl/Hermannstraße                                                         & NNO           & -                                                                                & \begin{tabular}[c]{@{}l@{}}10, 11, 12,\\ 13, 30\end{tabular} \\
10 & Brühl/Hermannstraße                                                         & OSO           & -                                                                                & 9, 11, 12                                                    \\
11 & Brühl/Hermannstraße                                                         & WNW           & -                                                                                & 9, 10, 12                                                    \\
12 & Brühl/Hermannstraße                                                         & SSW           & -                                                                                & 6, 9, 11                                                     \\
13 & Brühl                                                                       & NNO           & -                                                                                & 15, 16, 17                                                   \\
15 & Brühl                                                                       & SSW           & -                                                                                & 12, 14, 16                                                   \\
16 & Brühl                                                                       & WNW           & \begin{tabular}[c]{@{}l@{}}Rosa-Luxemburg-\\ Grundschule\\ am Brühl\end{tabular} & 14, 15                                                       \\
17 & Brühl/Elisenstraße                                                          & NNO           & Schriftzug \glqq Zuhause\grqq                                                             & 18, 19, 20, 21                                               \\
18 & Brühl/Elisenstraße                                                          & SSW           & -                                                                                & 15, 17, 19, 20                                               \\
19 & Brühl/Elisenstraße                                                          & OSO           & -                                                                                & 18, 20, 21                                                   \\
20 & Brühl/Elisenstraße                                                          & WNW           & -                                                                                & 18, 19, 21                                                   \\
21 & Brühl/Elisenstraße                                                          & NNO           & -                                                                                & 19, 20, 22, 27                                               \\
22 & Brühl/Elisenstraße                                                          & SSW           & -                                                                                & 18, 19, 20, 21                                               \\
24 & \begin{tabular}[c]{@{}l@{}}Brühl/Zöllnerstraße/\\ Zöllnerplatz\end{tabular} & SSW           & -                                                                                & 22, 25, 26, 27                                               \\
25 & \begin{tabular}[c]{@{}l@{}}Brühl/Zöllnerstraße/\\ Zöllnerplatz\end{tabular} & OSO           & -                                                                                & 24, 26, 27                                                   \\
26 & \begin{tabular}[c]{@{}l@{}}Brühl/Zöllnerstraße/\\ Zöllnerplatz\end{tabular} & WNW           & -                                                                                & 24, 25, 27, 28                                               \\
27 & \begin{tabular}[c]{@{}l@{}}Brühl/Zöllnerstraße/\\ Zöllnerplatz\end{tabular} & NNO           & -                                                                                & 24, 25, 26, 28                                               \\
28 & Zöllnerplatz                                                                & S             & Der verrückte Stuhl                                                              & 24                                                           \\
30 & Brühl/Hermannstraße                                                         & NW            & Urteil des Paris                                                                 & 9                                                           
\end{tabular}
\end{table}



\chapter{Implementierung}
Der Player ist in den grundlegenden Webtechniken HTML, CSS, JacaScript und XML umgesetzt wirden. Da keine serverseitigen Technologien wie etwa PHP oder ähnliche zum Einsatz kommen, ist kein Webserver notwendig, um auf den Player zugreifen zu können. Aus diesem Grund sind die Anforderungen an das System des Endgerätes, auf dem der Player betrachtet werden soll, sehr gering. Es wird lediglich eine funktionierende Internetverbindung, ein Videoausgabegerät, wie zum Beispiel ein Display, und eine Software, die die verwendeten Webtechnologien zu interpretieren weiß, benötigt.

\section{Verwendete Bibliotheken}
Folgende Bibliotheken wurden im Rahmen der Projektimplementierung genutzt:
\begin{itemize}
\item Video.js (Player) \\ Open-source HTML5 Video Player
\item Video.js-Overlay-Plugin (Player Overlays) \\ POI-Overlay
\item Bootstrap (UI Elements) \\ Grid-Layout (siehe Bootstrap-Dokumentation, z.B. col-md-4)
\item jQuery (common lib for JS) \\ Easy access and traversing through DOM Tree
\item XML2JSON \\ Convert XML to JSON-Format \footnote{\url{https://github.com/sparkbuzz/jQuery-xml2json}}
\end{itemize}
\pagebreak
\section{Aufbau der XML-Datei}
Die XML-Datei enthält den Wurzelknoten \glqq interactive\_media\grqq, in dem sich der Knoten \glqq videos\grqq{} befindet. Unter diesem Knoten sind jegliche Videos mit dem Tag \glqq video\grqq{} und einer jeweils einzigartigen Video-ID aufgelistet. Zu jedem Video sind folgende Informationen in der XML-Datei enthalten:
\begin{itemize}
\item Video-ID
\item Dateiname
\item Straße beziehungsweise Straßenkreuzung
\item Blickrichtung nach der Himmelsrichtung
\item Geographische Koordinaten der Position
\item Point-of-Interest mit Bezeichnung, Beschreibung und gegebenenfalls einem Link zur Website (falls vorhanden)
\item Verknüfte Nachbarvideos mit entsprechender Video-ID sowie der Art der Verknüpfung (Richtung)
\end{itemize}

\section{Installation}
Die Installation des Projektes erfolgt über das Klonen des entsprechenden Github-Repositories\footnote{\url{https://github.com/benkindt/video-bruehl}}, dem Überprüfen der funktionierenden Internetverbindung sowie dem anschließenden Öffnen des Players über die Datei \glqq video-bruehl/src/index.html\grqq. Weiterhin muss der Internetbrowser imstande sein, auf lokale Dateien zuzugreifen. Dies erfolgt beispielsweise im Falle das Google Chrome Browsers über den zusätzlichen Startparameter \glqq –allow-file-access-from-files\grqq.



\chapter{Zusammenfassung}


%======================================================================
%	Glossar
%======================================================================

%\manualmark
%\addcontentsline{toc}{chapter}{Glossar}
%\markboth{Glossar}{Glossar}
%\def\glossaryname{Glossar}
%\printglosstex(glo)
%\cleardoublepage

%======================================================================
%	Literaturverzeichnis
%======================================================================
\cleardoublepage
 \thispagestyle{empty}
\begin{flushleft}
\manualmark
\markboth{Literaturverzeichnis}{Literaturverzeichnis}
\bibliography{literatur}
\bibliographystyle{user}
\cleardoublepage
\end{flushleft}
%======================================================================
%	Thesen
%======================================================================
%\chapter{Thesen}
%\include{thesen}

%======================================================================
%	Selbstständigkeitserklärung
%======================================================================
\chapter*{Selbstständigkeitserklärung}

Hiermit erkläre ich, dass ich die vorliegende Arbeit
selbstständig angefertigt, nicht anderweitig zu Prüfungszwecken vorgelegt und
keine anderen als die angegebenen Hilfsmittel verwendet habe. Sämtliche 
wissentlich verwendete Textausschnitte, Zitate oder Inhalte anderer Verfasser 
wurden ausdrücklich als solche gekennzeichnet.\\[2ex]
\dcplace, den \dcdate\\[6ex]
\flushleft
\newlength\us
\settowidth{\us}{-\dcauthorfirstname~\dcauthorlastname-}
\begin{tabular}{p{\us}}\hline
\centering\footnotesize \dcauthorfirstname~\dcauthorlastname
\end{tabular}

%======================================================================
%	Anhang
%======================================================================


\part*{Anhang}
\cleardoublepage
\begin{appendix}
	
\include{B/Anhang}

\end{appendix}

\end{document}

